\documentclass[a4paper,11pt]{article}
\usepackage[margin=27mm]{geometry}
\usepackage{fontspec,amsmath,amssymb,amsthm,unicode-math,booktabs}
\setmainfont{STIXGeneral}\setmathfont{STIX Two Math}
\newcommand{\gen}[1]{\langle#1\rangle}
\newcommand{\mix}{\mathbin{\diamond}}\newcommand{\Aut}{\operatorname{Aut}}
\newcommand{\F}{\mathbb F_3}
\newtheorem{theorem}{Theorem}\newtheorem{lemma}[theorem]{Lemma}
\title{Two medial orbits require two embedded intersection tests}
\author{Human Compression project}
\date{Additive author edition; pending independent review}
\begin{document}\maketitle
\begin{abstract}
We construct a finite four-polytope in class $2_{\{1,2\}}$ with exactly two
orbits of medial sections. Its smallest regular cover is polytopal, but
the intersection condition printed in Cunningham--Hubard,
arXiv:2506.19334v1, Corollaries~6.10 and~6.15(2), is false.
The example has $144$ flags; its regular cover has $2592$ flags.
The distinction is between subgroups embedded at different paired
roots, rather than between abstract section types. The proof establishes
all polytopality and automorphism assertions. A complete finite replay
also checks the actual graphs and subgroup embeddings.
\end{abstract}

\section{A finite two-orbit polytope}
All coordinates lie in $\F$. Let $D$ be the affine group of the line
generated by $s(z)=-z$ and $t(z)=1-z$. Their product is a nonidentity
translation, so $D$ is dihedral of order six. On $\F^2$ put
\[
 a(x,y)=(-x,y),\quad b(x,y)=(x,1-y),\quad d(x,y)=(y,x),
 \quad a'=dad,\quad b'=dbd.                         \tag{1}
\]
Composition here is ordinary function composition.
The group $H=\gen{a,b,d}$ is $(D\times D)\rtimes\gen d$, of order
$72$: $a,a',b,b'$ generate the two copies of $D$, and $d$ interchanges
them.

Define a colored graph $P$ on the $144$ pairs $(h,\epsilon)$, with
$h\in H$ and $\epsilon\in\{0,1\}$, by
\[
\begin{array}{ll}
 r_0(h,\epsilon)=(h,1-\epsilon),&
 r_1(h,\epsilon)=(a_\epsilon h,\epsilon),\\
 r_2(h,\epsilon)=(bh,\epsilon),&
 r_3(h,\epsilon)=(dh,1-\epsilon),
\end{array}
\quad a_0=a,\quad a_1=a'.                           \tag{2}
\]
These are fixed-point-free involutions with distinct neighbors at each
flag. Nonadjacent colors commute; the only case depending on (1) is
$a_{1-\epsilon}d=da_\epsilon$. The graph is connected: colors $0,3$
supply multiplication by $d$ without changing $\epsilon$, while at
$\epsilon=0$ the other colors supply $a,b$. Thus (2) is a maniplex.

We use the standard path-intersection characterization of polytopality:
at every flag, orbits under two color sets intersect in the orbit under
their intersection. Its recursive form says that polytopality of all
facets and vertex-figures, together with their exact intersection at each
flag, suffices. The companion common paper proves this recursive form
from the established maniplex--polytope correspondence. It applies to
digons as well as to polygons with more sides.

\begin{lemma}\label{input}
The graph $P$ is a four-polytope. Its full color-preserving automorphism
group is the right regular action of $H$. It has class $2_{\{1,2\}}$,
exactly two medial orbits, and one orbit each of facets and vertex-figures.
Its medial sections are digons and triangles.
\end{lemma}
\begin{proof}
Right multiplication $R_k(h,\epsilon)=(hk,\epsilon)$ commutes with
all adjacencies and is transitive on each set defined by $\epsilon$.
Thus it suffices to check local intersections at $(1,0)$ and $(1,1)$.
Products below refer to the two coordinate actions on $\F^2$.
The facet through either base flag has the same $H$-coordinates in both
sets, namely $A_H=\gen a\times D$.
The vertex-figure through $(1,0)$ has coordinates
$B_H=D\times\gen b$ at $\epsilon=0$ and $dB_H$ at $\epsilon=1$.
Through $(1,1)$ these are $B'_H=\gen{b'}\times D$ at $\epsilon=1$
and $dB'_H$ at $\epsilon=0$. Eliminating the color which changes
$\epsilon$ proves these descriptions; for example the even words in
colors $1,2,3$ generate $a,b,b'$ in the first case.
Diagonal coordinate actions and actions which interchange coordinates
are disjoint. Moreover,
\[
 A_H\cap B_H=\gen{a,b},\qquad
 A_H\cap B'_H=\gen{a',b}.                          \tag{3}
\]
The distinct reflections $s,t$ have trivial subgroup intersection on
the first coordinate. The two right sides in (3) have orders four and
six, respectively. This gives the required top intersection at both flags.

For the facet check, the even part of the colors $0,1$ orbit is
$\gen{a,a'}$, and the odd part has the same coordinates.
Its intersections with the two medial groups are $\gen a$ and
$\gen{a'}$, respectively. For the vertex-figure check, the even part
of the colors $2,3$ orbit is $\gen{b,b'}$, and the odd part is its
left translate by $d$. Its intersection with either medial group is
$\gen b$. These assertions again use $\gen s\cap\gen t=1$.
All rank-two residues are alternating cycles and therefore polygons.
The recursive criterion proves both kinds of rank-three residues and
then, by (3), the whole polytope.

An automorphism of a connected colored graph is determined by the
image of one flag. One interchanging the two sets defined by $\epsilon$
would send a middle-color component of size four to one of size six,
which is impossible. The right action of $H$ is already regular on each
set, so it is the full automorphism group. Colors $1,2$ preserve its
orbits; colors $0,3$ interchange them. Every medial component stays in
one set, and $H$ is transitive there: the two different sizes give
precisely two medial orbits. Each facet and vertex-figure meets both
sets, so $H$ is transitive on those faces. In particular all facets are
isomorphic and all vertex-figures are isomorphic, as the source requires.
\end{proof}

\section{The regular cover and its fibers}
Let $\overline P$ denote $P$ rooted at $(1,1)$ instead of $(1,0)$.
The connected mix $Q=P\mix\overline P$ starts at
\[
 q_0=((1,0),(1,1)),\qquad q_1=r_0q_0=((1,1),(1,0)).
\]
Write its flags as $(h_0,h_1,\epsilon)$, representing
$((h_0,\epsilon),(h_1,1-\epsilon))$. Those with $\epsilon=0$ are
exactly the pairs in
\[
 K=\gen{(a,a'),(b,b),(d,d)}\leq H\times H.           \tag{4}
\]
The displayed generators are realized by colors $1$, $2$, and $0,3$.
Conversely, eliminate color $0$ from an even word; its conjugation sends
$(a,a')$ to $(a',a)$, already the conjugate by $(d,d)$, and fixes the
other two generators. Flags at $\epsilon=1$ have the same pairs, since
color $0$ changes only $\epsilon$. Right multiplication by $K$ acts
regularly on each set of flags. Interchanging the two input flags sends
$q_0$ to $q_1$, preserves the connected component, and commutes with
all colors. Consequently $Q$ is regular.

For completeness $|K|=1296$. The squares of products of $(b,b)$
with $(a,a')$ and with $(a',a)$ give independent translations in the
two $y$-coordinates. Conjugation by $(d,d)$ gives both $x$-translations.
Hence $K$ contains all $3^4$ translations. Its linear image has the
three independent sign vectors
\[
 (1,0,0,1),\quad(0,1,1,0),\quad(0,1,0,1)
\]
and simultaneous interchange of the two axes in both factors.
The coordinate order here is $x_0,y_0,x_1,y_1$; a $1$ denotes a sign
change. The span is stable under the interchange, which has order two
outside the sign subgroup. Thus the linear image has order $16$,
and $Q$ has $2\cdot3^4\cdot16=2592$ flags.

For a color set $J$ and $q_k=(x_k,y_k)$ define the actual fiber
\[
 F_J(q_k)=\{h\in H:(x_k,R_hy_k)\text{ lies in the }J
                  \text{-component of }q_k\}.
\]
These are the embedded lower variance subgroups in the second input's
automorphism group. The right regular action identifies this group with
$H^{\mathrm{op}}$, which changes none of its subgroups or intersections.
Set $A_k=F_{\{0,1,2\}}(q_k)$, $B_k=F_{\{1,2,3\}}(q_k)$ and
$C_k=F_{\{1,2\}}(q_k)$. For $u\in\{0,1\}$ write $\eta=(-1)^u$.
Inside this single specified copy of $\Aut(P)$ we have
\[
\begin{array}{ll}
 A_0=A_1=\{(x,y)\mapsto(\eta x,\eta y+c):u=0,1,\ c\in\F\},\\
 B_0=\{(x,y)\mapsto(\eta x+u,\eta y+c):u=0,1,\ c\in\F\},\\
 B_1=\{(x,y)\mapsto(\eta x+c,\eta y+u):u=0,1,\ c\in\F\},\\
 C_0=\{(x,y)\mapsto(x,y+c):c\in\F\},\qquad
 C_1=\gen{(x,y)\mapsto(-x,1-y)}.
\end{array}                                                       \tag{5}
\]
Here are complete generation checks. Put $U=(a,a')$, $W=(a',a)$,
$V=(b,b)$ and $V'=(b',b')$. The even part of a mixed facet is
$\gen{U,W,V}$, consisting exactly of the pairs
\[
\begin{split}
 h_0(x,y)&=((-1)^u x,\ (-1)^{w+v}y+\alpha),\\
 h_1(x,y)&=((-1)^w x,\ (-1)^{u+v}y+\beta),
\end{split}                                                       \tag{6}
\]
where $u,w,v\in\{0,1\}$ and $\alpha,\beta\in\F$.
The generators give the independent signs, while $(UV)^2$ and
$(WV)^2$ give separate nonzero translations on the two $y$-coordinates.
Thus every displayed pair occurs. Setting $h_0=1$ gives $A_0$;
factor interchange gives the same set for $A_1$.

The even part of the mixed vertex-figure at $q_0$ is
$\gen{U,V,V'}$. Its elements are precisely
\[
\begin{split}
 h_0(x,y)&=((-1)^{u+w}x+\alpha,\ (-1)^v y+v),\\
 h_1(x,y)&=((-1)^w x+w,\ (-1)^{u+v}y+\beta).
\end{split}                                                       \tag{7}
\]
The signs are independent. The squares $(UV)^2$ and $(UV')^2$ give
independent translations on $y_1$ and $x_0$. For each remaining
coordinate the map is $z$ or $1-z$, so its translation is forced by
its sign; these conditions are closed under composition. This proves
equality in (7). Setting $h_0=1$ yields $B_0$, and setting $h_1=1$
yields $B_1$ at the reversed paired root.
Finally the middle group $\gen{U,V}$ is dihedral of order twelve:
its product has order six, with first projection $ab$ of order two
and second projection $a'b$ of order three. Its projection kernels
have orders three and two and give exactly the two groups $C_0,C_1$.
Consequently
\[
 A_0\cap B_0=C_0\cong C_3,\qquad
 A_1\cap B_1=C_1\cong C_2,\qquad C_0\cap C_1=1.       \tag{8}
\]

\begin{theorem}\label{counterexample}
The smallest regular cover of $P$ is polytopal. Nevertheless
$P\mix\overline P$ violates the condition printed in
[CH, Corollary~6.10(3)] and the medial variance condition printed in
[CH, Corollary~6.15(2)]. All exceptional two-orbit hypotheses hold.
\end{theorem}
\begin{proof}
Every facet and vertex-figure of $Q$ is a mix of two polyhedra and is
polytopal. Here the usual polyhedron-mix result also has a short path
proof. An intersection of the two rank-three end orbits has its
second-coordinate fiber inside one middle edge of the second
polyhedron. Its nonidentity endpoint would require a word fixing a
polygon flag in the first factor but changing polygon flag parity
in the second. Polygon flag graphs are bipartite, so this is
impossible. The common fiber-intersection lemma proves the rank-three
intersection property, including when a polygon is a digon.

At $q_0$, (8) is exactly the top fiber-intersection condition.
Regularity of $Q$ transports the resulting path condition to every
flag, whether or not that automorphism projects to automorphisms of
the individual rooted factors. The recursive criterion therefore
proves polytopality of $Q$.

Let $W$ be the universal string connection group and $N$ the stabilizer
of $(1,0)$ in $P$. Its index-two subgroup $V$ preserving $\epsilon$
normalizes $N$: on each of its two orbits the right group $H$ acts
regularly and centralizes $V$, so the image of $V$ acts regularly as well.
Thus the core of $N$ in $W$ is $N\cap r_0Nr_0$, the stabilizer of
$q_0$ in the mix. It follows that $Q$ is the smallest regular
premaniplex cover. Since it is polytopal, it is also the smallest
regular polytopal cover.

Lemma~\ref{input} shows that both inputs genuinely belong to
$2_{\{1,2\}}$, with uniform facet and vertex-figure types and precisely
two medial orbits. Corollary~6.10(3) would require
$A_0\cap B_0\leq C_0\cap C_1$. Formula (8) contradicts it:
the left side has order three and the right side is trivial.
A literal witness is the color sequence $1,2,1,2$ at $q_0$.
It fixes the first flag and moves the second by
$(x,y)\mapsto(x,y+2)$, a nonidentity member of both end fibers outside
$C_1$.

For Corollary~6.15(2), the two medial sections are a digon and a
triangle. Their two lower variance groups have orders three and two,
so their intersection is trivial under any embeddings in one ambient
group. The printed order $X(K\mid\overline K)$ puts $P$ second,
so use the reversed factor order $\overline P\mix P$: in our single
ambient copy of $H$ this is the calculation at $q_1$, with common
end variance group $A_1\cap B_1=C_1$ of order two. Both end mixes
are polytopal, but this nontrivial group cannot lie in the stated
trivial intersection. Reversing the initial base flag gives instead
order three on the left, with the same failure. No regular-input
degeneration or identification of abstract group types is being used.
\end{proof}

\section{Corrected criterion and scope}
The valid exceptional criterion keeps the two actual paired roots:
\[
 A_0\cap B_0\leq C_0,\qquad A_0\cap B_1\leq C_1.       \tag{9}
\]
The equality $A_0=A_1$ holds because the flags lie in one mixed facet.
Their vertex-figures are isomorphic but are different faces, giving
no equality of the embedded stabilizers $B_0,B_1$.
Even the weaker proposed equality $A_0\cap B_0=A_0\cap B_1$ fails,
by the different orders in (8). The companion paper proves (9) in
general. For the smallest regular cover, regularity permits one
correctly rooted test; it does not permit intersecting two differently
transported medial subgroups on its right.

This addition resolves the two qualified entries in our frozen F005
recovery packet without changing that packet. It does not answer the
three final problems of [CH] or assert a compression factor or
historical priority. The universal connection-group, path-intersection
and regular-cover correspondences are retained standard inputs. The
affine construction and all the intersections above are proved here.
The complete program is a validation of these written proofs, not a
premise replacing a classification or geometric argument.

\begin{thebibliography}{9}
\bibitem{CH} G.~Cunningham and I.~Hubard, \emph{Polytopality criteria for the
mix of polytopes and maniplexes}, arXiv:2506.19334v1, 24 June 2025,
physical pp.~17 and 19, Corollaries~6.10 and 6.15(2).
\bibitem{fiber} Human Compression project, \emph{Intersections in simultaneous group
actions} and \emph{Polytopality of mixes and their regular covers},
frozen author packet \texttt{topic\_Mix\_F005\_recovery}, 2026;
rooted exceptional tests and their exact common-facet equality.
\end{thebibliography}
\end{document}
